\section{1.ª ley modular para variaciones del àrea}

Definamos las dos \`areas normalizadas
\begin{equation}
\label{rm:eq:modular-normalized-area-channels}
\mathcal A_m=\gammaH a_0N_m,
\qquad m\in\{90,120\}.
\end{equation}
Entonces
\begin{equation}
\label{rm:eq:modular-beta-definition}
\HHM^{(A)}-cI
=\beta_{90}\mathcal A_{90}+\beta_{120}\mathcal A_{120},
\qquad
\beta_m=\frac{m\Dmod}{\gammaH a_0}.
\end{equation}
Los valores terminales son
\begin{align}
\beta_{90}&=5.9726485401070929914651798889\ldots,
\label{rm:eq:modular-beta90}\\
\beta_{120}&=7.9635313868094573219535731852\ldots,
\qquad
\frac{\beta_{120}}{\beta_{90}}=\frac43.
\label{rm:eq:modular-beta120}
\end{align}

\begin{theorem}[Primera ley modular en variables areales]
\label{rm:thm:modular-first-law}
Para una familia diferenciable de estados normalizados
\(\rho(\varepsilon)\) que pasa por \(\rho_A\), manteniendo fijos
\(\Dmod,\gammaH,a_0\),
\begin{equation}
\label{rm:eq:modular-first-law}
\boxed{
\delta S
=\beta_{90}\,\delta\langle\mathcal A_{90}\rangle
+\beta_{120}\,\delta\langle\mathcal A_{120}\rangle.}
\end{equation}
\end{theorem}

\begin{proof}
La primera ley de entrelazamiento en el estado de referencia es
\(\delta S=\delta\langle-\log\rho_A\rangle\). El t\'ermino escalar
\(cI\) no contribuye porque \(\delta\Tr\rho=0\). Sustituir
\cref{rm:eq:modular-beta-definition} prueba la identidad.
\end{proof}

La f\'ormula es local en el espacio de estados. Para cambios finitos aparece
la entrop\'ia relativa; omitirla convertir\'ia una identidad tangente en una
afirmaci\'on falsa de linealidad global.

\Needspace{22\baselineskip}
\begin{theorem}[Identidad finita con t\'ermino de defecto]
\label{rm:thm:modular-finite-first-law}
Para todo estado \(\sigma\) cuyo soporte est\'e contenido en el de
\(\rho_A\),
\begin{equation}
\label{rm:eq:modular-finite-first-law}
\boxed{
S(\sigma)-S(\rho_A)
=\beta_{90}\Delta_\sigma\langle\mathcal A_{90}\rangle
+\beta_{120}\Delta_\sigma\langle\mathcal A_{120}\rangle
-D(\sigma\Vert\rho_A),}
\end{equation}
donde
\(D(\sigma\Vert\rho_A)=\Tr\sigma(\log\sigma-\log\rho_A)\geq0\).
En particular, la variaci\'on lineal del
\cref{rm:thm:modular-first-law} es el t\'ermino tangente de esta identidad.
\end{theorem}

\begin{proof}
Por definici\'on y por \(\HHM^{(A)}=-\log\rho_A\),
\[
D(\sigma\Vert\rho_A)
=-S(\sigma)+\Tr(\sigma\HHM^{(A)}).
\]
Para \(\sigma=\rho_A\), el miembro izquierdo es cero y
\(S(\rho_A)=\Tr(\rho_A\HHM^{(A)})\). Restar ambas identidades da
\[
S(\sigma)-S(\rho_A)
=\Delta_\sigma\langle\HHM^{(A)}\rangle
-D(\sigma\Vert\rho_A).
\]
El escalar \(cI\) se cancela entre estados normalizados; sustituir
\cref{rm:eq:modular-beta-definition} produce
\cref{rm:eq:modular-finite-first-law}. La positividad de la entrop\'ia
relativa es la desigualdad de Klein.
\end{proof}

La correcci\'on \(D(\sigma\Vert\rho_A)\) mide la distancia informacional
entre el estado real y el estado modular de referencia. As\'i, la primera
ley no pierde contenido al salir del r\'egimen infinitesimal: se convierte
en una identidad exacta con un defecto positivo controlado.
